\input{../tex-inputs/latex-leseansicht-vorspann}

\section[Arthur und Olga Schnitzler an Gustav Schwarzkopf, {[}13.?{]} 5. 1912]{L04454 Arthur und Olga Schnitzler an Gustav Schwarzkopf, {[}13.?{]} 5. 1912}
\nopagebreak\mylabel{L04454v}
\rehead{ }\normalsize\beginnumbering\briefempfaengerindex{Schwarzkopf, Gustav@\textsc{Schwarzkopf, Gustav}!zzzSchnitzler, Olga@\emph{von Olga Schnitzler}!1912-05-132@{{[}13.?{]} 5. 1912}|(be}\briefempfaengerindex{Schwarzkopf, Gustav@\textsc{Schwarzkopf, Gustav}!zzzSchnitzler, Arthur@\emph{von Arthur Schnitzler}!1912-05-132@{{[}13.?{]} 5. 1912}|(be}
\toendnotes[C]{\smallbreak\pagebreak[2]}
\correspDesc{Versand  durch Arthur Schnitzler, Olga Schnitzler am [13.?] 5. 1912 in Brijuni
\newline{}Übermittlung  am 14. 5. 1912 in Triest
\newline{}Erhalt  durch Gustav Schwarzkopf im Zeitraum [14. 5. 1912
                  – 18. 5. 1912?] in Wien}\toendnotes[C]{\smallbreak}
\Standort{DLA, A:Schnitzler, HS.1985.1.1897.}
\physDesc{Bildpostkarte, 160 Zeichen
\newline{}Handschrift Arthur Schnitzler: Bleistift, deutsche Kurrent
\newline{}Handschrift Olga Schnitzler: Bleistift, lateinische Kurrent
\newline{}Versand: Stempel: »\nobreak{}\oindex{Triest@\textbf{Triest}|pwk}Triest, 14. V. 12, 4, Trieste\nobreak{}«.  }\toendnotes[C]{\smallbreak}\pstart{}{\pb}Herrn Guſtav Schwarzkopf\pend{}\pstart{}Wien \substVorne{}\textsuperscript{X}\substDazwischen{}I\substHinten{}\oindex{I., Innere Stadt@\textbf{I., Innere Stadt}|pw}\pend{}\pstart{}\textsc{Tiefer Graben 17}\oindex{Wien@\textbf{Wien}!I., Innere Stadt@\textbf{I., Innere Stadt}!Tiefer Graben 17@\textbf{Tiefer Graben 17}|pw}\pend{}{\bigskip}
\pstart
           \noindent{}\centering{}{\pb}\textcolor{gray}{\textbf{Brioni-Paradies\oindex{Monte Paradiso@\textbf{Monte Paradiso}|pw}.}}\pend
           
\pstart
           \centering{}{[}hs. A. Schnitzler:{]} \label{T_L04454-1v}\edtext{Brioni\oindex{Brijuni@\textbf{Brijuni}|pw}}{\lemma{\textnormal{\emph{Brioni}}}\Cendnote{\textnormal{Schnitzler beschriftet das Bildmotiv in der
                  Mitte der Bildseite.}}}\label{T_L04454-1}\pend
           \vspace{1em}
\pstart
           \noindent{}{\pb}Herzliche Grüße. Es iſt \label{K_L04454-1v}\edtext{hier}{\lemma{\textnormal{\emph{hier}}}\Cendnote{\textnormal{Obwohl die
                     von den Schnitzlers\pwindex{Schnitzler, Olga 17.\,1.\,1882 Wien – 13.\,1.\,1970 Lugano@\textsc{Schnitzler, Olga} (17.\,1.\,1882 Wien – 13.\,1.\,1970 Lugano)|pwk} nicht datierte Karte laut Stempel am 14. 5. 1912 in Triest\oindex{Triest@\textbf{Triest}|pwk} aufgegeben wurde, ist sie augenscheinlich noch in Brioni\oindex{Brijuni@\textbf{Brijuni}|pwk} verfasst
                  worden und dürfte somit bereits am 13. 5. 1912 entstanden sein.}}}\label{K_L04454-1} über \textcolor{gray}{alles
                  Erwarten} { }ſchön. Wir fahren nach Trieſt\oindex{Triest@\textbf{Triest}|pw} – Venedig\oindex{Venedig@\textbf{Venedig}|pw}.\pend
           \pstart Ihr \spacefill\mbox{Arthur}\pend{}\selectlanguage{ngerman}\vspace{1em}
\pstart
           \noindent{}Herzlichst {\\}\spacefill\mbox{Olga.}\pend
           \selectlanguage{ngerman}\endnumbering\briefempfaengerindex{Schwarzkopf, Gustav@\textsc{Schwarzkopf, Gustav}!zzzSchnitzler, Olga@\emph{von Olga Schnitzler}!1912-05-132@{{[}13.?{]} 5. 1912}|)be}\briefempfaengerindex{Schwarzkopf, Gustav@\textsc{Schwarzkopf, Gustav}!zzzSchnitzler, Arthur@\emph{von Arthur Schnitzler}!1912-05-132@{{[}13.?{]} 5. 1912}|)be}\mylabel{L04454h}
\begin{anhang}
\end{anhang}\newcommand{\dateiname}{L04454}\newcommand{\titel}{Arthur und Olga Schnitzler an Gustav Schwarzkopf, [13.?] 5. 1912}\newcommand{\editorInnen}{Herausgegeben von Jahnke, SelmaMüller, Martin Anton}\input{../tex-inputs/latex-leseansicht-abspann}
